% Archivo: simbolos.tex — Lista de Símbolos

\newglossaryentry{alpha}{
  type=simbolos,
  name={$\alpha$},
  description={Ángulo de ataque},
  unit={°},
  sort={alpha}
}

\newglossaryentry{alphaz}{
  type=simbolos,
  name={$\alpha_z$},
  description={Exponente de cizalladura vertical del perfil de viento normal (NWP); no confundir con el ángulo de ataque $\alpha$},
  unit={---},
  sort={alpha-z}
}

\newglossaryentry{area}{
  type=simbolos,
  name={$A$},
  description={Área barrida por el rotor},
  unit={m$^2$},
  sort={A}
}

\newglossaryentry{cuerda}{
  type=simbolos,
  name={$c$},
  description={Cuerda de la pala del rotor; también, parámetro de escala de la distribución de Weibull},
  unit={m},
  sort={c}
}

\newglossaryentry{LM}{
  type=simbolos,
  name={$LM$},
  description={Multiplicador de carga del análisis de pandeo lineal; factor por el que deben amplificarse las cargas aplicadas para alcanzar la inestabilidad elástica},
  unit={---},
  sort={LM}
}

\newglossaryentry{cETM}{
  type=simbolos,
  name={$c_{ETM}$},
  description={Constante del modelo de turbulencia extrema (ETM)},
  unit={m/s},
  sort={cETM}
}

\newglossaryentry{Cd90}{
  type=simbolos,
  name={$C_{d90}$},
  description={Coeficiente de arrastre a 90°},
  unit={---},
  sort={Cd90}
}

\newglossaryentry{Cm}{
  type=simbolos,
  name={$C_m$},
  description={Coeficiente de momento},
  unit={---},
  sort={Cm}
}

\newglossaryentry{Cp}{
  type=simbolos,
  name={$C_p$},
  description={Coeficiente de potencia},
  unit={---},
  sort={Cp}
}

\newglossaryentry{altura}{
  type=simbolos,
  name={$H$},
  description={Altura del rotor},
  unit={m},
  sort={H}
}

\newglossaryentry{lambda}{
  type=simbolos,
  name={$\lambda$},
  description={Relación de velocidad de punta (TSR)},
  unit={---},
  sort={lambda}
}

\newglossaryentry{mu}{
  type=simbolos,
  name={$\mu$},
  description={Viscosidad dinámica del fluido},
  unit={Pa$\cdot$s},
  sort={mu}
}

\newglossaryentry{nu}{
  type=simbolos,
  name={$\nu$},
  description={Viscosidad cinemática del fluido; también, coeficiente de Poisson del material},
  unit={m$^2$/s},
  sort={nu}
}

\newglossaryentry{omega}{
  type=simbolos,
  name={$\Omega$},
  description={Velocidad angular del rotor},
  unit={rad/s},
  sort={Omega}
}

\newglossaryentry{theta}{
  type=simbolos,
  name={$\theta$},
  description={Posición acimutal de la pala},
  unit={°},
  sort={theta}
}

\newglossaryentry{radio}{
  type=simbolos,
  name={$R$},
  description={Radio del rotor},
  unit={m},
  sort={R}
}

\newglossaryentry{reynolds}{
  type=simbolos,
  name={$Re$},
  description={Número de Reynolds},
  unit={---},
  sort={Re}
}

\newglossaryentry{rho}{
  type=simbolos,
  name={$\rho$},
  description={Densidad del fluido},
  unit={kg/m$^3$},
  sort={rho}
}

\newglossaryentry{Vave}{
  type=simbolos,
  name={$V_{ave}$},
  description={Velocidad media anual del viento a la altura del centro del rotor},
  unit={m/s},
  sort={Vave}
}

\newglossaryentry{Vhub}{
  type=simbolos,
  name={$V_{hub}$},
  description={Velocidad del viento a la altura del centro del rotor},
  unit={m/s},
  sort={Vhub}
}

\newglossaryentry{zhub}{
  type=simbolos,
  name={$z_{hub}$},
  description={Altura del centro del rotor},
  unit={m},
  sort={zhub}
}

% ---- Geometría y parámetros del rotor ----
\newglossaryentry{D}{
  type=simbolos,
  name={$D$},
  description={Diámetro del rotor},
  unit={m},
  sort={D}
}

\newglossaryentry{Npalas}{
  type=simbolos,
  name={$N$},
  description={Número de palas},
  unit={---},
  sort={N-palas}
}

\newglossaryentry{sigma}{
  type=simbolos,
  name={$\sigma$},
  description={Solidez del rotor ($\sigma = Nc/R$). La tensión mecánica se escribe $\sigma(t)$ cuando varía en el tiempo, y con subíndice en los demás casos},
  unit={---},
  sort={sigma}
}

\newglossaryentry{DeltaTheta}{
  type=simbolos,
  name={$\Delta\theta$},
  description={Paso angular de discretización temporal},
  unit={°},
  sort={DeltaTheta}
}

% ---- Aerodinámica ----
\newglossaryentry{CL}{
  type=simbolos,
  name={$C_l$},
  description={Coeficiente de sustentación},
  unit={---},
  sort={CL}
}

\newglossaryentry{CD}{
  type=simbolos,
  name={$C_d$},
  description={Coeficiente de arrastre},
  unit={---},
  sort={CD}
}

\newglossaryentry{M}{
  type=simbolos,
  name={$M$},
  description={Momento aerodinámico (torque)},
  unit={N$\cdot$m},
  sort={M}
}

\newglossaryentry{P}{
  type=simbolos,
  name={$P$},
  description={Potencia mecánica del rotor},
  unit={W},
  sort={P}
}

\newglossaryentry{W}{
  type=simbolos,
  name={$W$},
  description={Velocidad relativa del flujo sobre la pala},
  unit={m/s},
  sort={W}
}

\newglossaryentry{Ncrit}{
  type=simbolos,
  name={$N_{crit}$},
  description={Parámetro de transición del modelo $e^N$},
  unit={---},
  sort={Ncrit}
}

% ---- Recurso eólico y oleaje ----
\newglossaryentry{Vref}{
  type=simbolos,
  name={$V_{ref}$},
  description={Velocidad de referencia del viento (período de retorno de 50 años)},
  unit={m/s},
  sort={Vref}
}

\newglossaryentry{z0}{
  type=simbolos,
  name={$z_0$},
  description={Longitud de rugosidad de la superficie del mar},
  unit={m},
  sort={z0}
}

\newglossaryentry{AC}{
  type=simbolos,
  name={$A_C$},
  description={Constante de Charnock; la IEC 61400-3 la fija en $0{,}011$ para mar abierto y $0{,}034$ para emplazamientos cercanos a la costa},
  unit={---},
  sort={AC}
}

\newglossaryentry{kappa}{
  type=simbolos,
  name={$\kappa$},
  description={Constante de von K\'arm\'an ($0{,}4$)},
  unit={---},
  sort={kappa}
}

\newglossaryentry{Iref}{
  type=simbolos,
  name={$I_{ref}$},
  description={Intensidad de turbulencia de referencia, evaluada a $V_{hub} = 15$~m/s, según la categoría de turbina (A, B o C); sin relación con $V_{ref}$},
  unit={---},
  sort={Iref}
}

\newglossaryentry{bNTM}{
  type=simbolos,
  name={$b_{NTM}$},
  description={Parámetro de ordenada del Modelo de Turbulencia Normal (NTM); la IEC 61400-1 lo denota $b$ y lo fija en $5{,}6$~m/s para las clases estándar},
  unit={m/s},
  sort={b-NTM}
}

\newglossaryentry{sigma1}{
  type=simbolos,
  name={$\sigma_{NTM}$},
  description={Desviación estándar representativa de la turbulencia según el Modelo de Turbulencia Normal (NTM); la IEC 61400-1 la denota $\sigma_1$},
  unit={m/s},
  sort={sigma-NTM}
}

\newglossaryentry{fWeibull}{
  type=simbolos,
  name={$f_W$},
  description={Densidad de probabilidad de Weibull de la velocidad del viento en el emplazamiento},
  unit={s/m},
  sort={fWeibull}
}

\newglossaryentry{kWeibull}{
  type=simbolos,
  name={$k$},
  description={Parámetro de forma de la distribución de Weibull},
  unit={---},
  sort={kWeibull}
}

% ---- Material de la pala ----
\newglossaryentry{Emod}{
  type=simbolos,
  name={$E$},
  description={Módulo elástico del material},
  unit={GPa},
  sort={Emod}
}

\newglossaryentry{sigmau}{
  type=simbolos,
  name={$\sigma_u$},
  description={Resistencia última del material},
  unit={MPa},
  sort={sigma-u}
}

\newglossaryentry{ml}{
  type=simbolos,
  name={$m_l$},
  description={Masa por unidad de longitud de la pala},
  unit={kg/m},
  sort={ml}
}

\newglossaryentry{teq}{
  type=simbolos,
  name={$t_{eq}$},
  description={Espesor equivalente de pared de la pala},
  unit={mm},
  sort={teq}
}

% ---- Modelo de turbulencia de Mann ----
\newglossaryentry{epsilonmann}{
  type=simbolos,
  name={$\varepsilon$},
  description={Tasa de disipación turbulenta; interviene en el modelo de Mann a través del parámetro $\alpha\varepsilon^{2/3}$. No confundir con la deformación última $\varepsilon_u$},
  unit={m$^2$/s$^3$},
  sort={epsilon}
}

\newglossaryentry{epsilonu}{
  type=simbolos,
  name={$\varepsilon_u$},
  description={Deformación última del material},
  unit={\%},
  sort={epsilon-u}
}

\newglossaryentry{sigmaiso}{
  type=simbolos,
  name={$\sigma_{iso}$},
  description={Parámetro de intensidad del modelo de Mann, $\sigma_{iso} = 0{,}55\,\sigma_{NTM}$},
  unit={m/s},
  sort={sigma-iso}
}

\newglossaryentry{Lambda1}{
  type=simbolos,
  name={$\Lambda_1$},
  description={Parámetro de escala de turbulencia longitudinal (IEC 61400-1, Sección 6.3)},
  unit={m},
  sort={Lambda1}
}

\newglossaryentry{ellmann}{
  type=simbolos,
  name={$\ell$},
  description={Escala de longitud del modelo de Mann, $\ell = 0{,}8\,\Lambda_1$},
  unit={m},
  sort={ell}
}

\newglossaryentry{Gammamann}{
  type=simbolos,
  name={$\Gamma$},
  description={Parámetro de anisotropía del modelo de Mann ($\Gamma = 3{,}9$ según IEC 61400-1, donde se denota $\gamma$)},
  unit={---},
  sort={Gamma}
}

\newglossaryentry{gammaf}{
  type=simbolos,
  name={$\gamma_f$},
  description={Factor parcial de seguridad aplicado a las cargas},
  unit={---},
  sort={gamma-f}
}

\newglossaryentry{gammam}{
  type=simbolos,
  name={$\gamma_m$},
  description={Factor parcial de seguridad del material},
  unit={---},
  sort={gamma-m}
}

\newglossaryentry{gamman}{
  type=simbolos,
  name={$\gamma_n$},
  description={Factor parcial de seguridad por consecuencias de la falla},
  unit={---},
  sort={gamma-n}
}

\newglossaryentry{gammaE}{
  type=simbolos,
  name={$\gamma_E$},
  description={Constante de Euler-Mascheroni ($0{,}5772$), empleada en el ajuste de Gumbel},
  unit={---},
  sort={gamma-E}
}

\newglossaryentry{bbasquin}{
  type=simbolos,
  name={$b$},
  description={Exponente de Basquin de la curva S-N del material},
  unit={---},
  sort={b}
}

\newglossaryentry{Ftot}{
  type=simbolos,
  name={$F_{tot}$},
  description={Fuerza total integrada sobre la pala},
  unit={N},
  sort={F-tot}
}

\newglossaryentry{Fmax}{
  type=simbolos,
  name={$F_{max}$},
  description={Máximo de $F_{tot}(t)$ en una realización},
  unit={N},
  sort={F-max}
}

\newglossaryentry{Fk}{
  type=simbolos,
  name={$F_k$},
  description={Carga característica obtenida de la extrapolación estadística},
  unit={N},
  sort={F-k}
}

\newglossaryentry{Fd}{
  type=simbolos,
  name={$F_d$},
  description={Carga de diseño, $F_d = \gamma_f F_k$},
  unit={N},
  sort={F-d}
}

\newglossaryentry{Pe}{
  type=simbolos,
  name={$P_e$},
  description={Probabilidad de excedencia},
  unit={---},
  sort={P-e}
}

\newglossaryentry{Tobs}{
  type=simbolos,
  name={$T_{obs}$},
  description={Duración de cada realización simulada},
  unit={s},
  sort={T-obs}
}

\newglossaryentry{Tretorno}{
  type=simbolos,
  name={$T_{retorno}$},
  description={Período de retorno; la IEC 61400-1 lo fija en $50$ años para el DLC 1.1},
  unit={años},
  sort={T-retorno}
}

\newglossaryentry{Tsim}{
  type=simbolos,
  name={$T_{sim}$},
  description={Duración efectiva de la serie temporal procesada},
  unit={s},
  sort={T-sim}
}

\newglossaryentry{muj}{
  type=simbolos,
  name={$\mu_j$},
  description={Parámetro de localización de la distribución de Gumbel a la velocidad $V_j$},
  unit={N},
  sort={mu-j}
}

\newglossaryentry{betaj}{
  type=simbolos,
  name={$\beta_j$},
  description={Parámetro de escala de la distribución de Gumbel a la velocidad $V_j$},
  unit={N},
  sort={beta-j}
}

\newglossaryentry{sigma1max}{
  type=simbolos,
  name={$\sigma_{1,\max}$},
  description={Tensión principal máxima obtenida del modelo de elementos finitos},
  unit={MPa},
  sort={sigma-1max}
}

\newglossaryentry{sigma3min}{
  type=simbolos,
  name={$\sigma_{3,\min}$},
  description={Tensión principal mínima obtenida del modelo de elementos finitos},
  unit={MPa},
  sort={sigma-3min}
}

\newglossaryentry{sigmaa}{
  type=simbolos,
  name={$\sigma_a$},
  description={Amplitud de tensión de un ciclo de fatiga},
  unit={MPa},
  sort={sigma-a}
}

\newglossaryentry{sigmam}{
  type=simbolos,
  name={$\sigma_m$},
  description={Tensión media de un ciclo de fatiga},
  unit={MPa},
  sort={sigma-m}
}

\newglossaryentry{sigmain}{
  type=simbolos,
  name={$\sigma_{in}$},
  description={Tensión principal máxima en la sección crítica bajo cargas inerciales},
  unit={MPa},
  sort={sigma-in}
}

\newglossaryentry{Keff}{
  type=simbolos,
  name={$K_{aero}$},
  description={Factor que relaciona el momento flector aerodinámico con la tensión principal},
  unit={1/m$^3$},
  sort={K-aero}
}

\newglossaryentry{Mres}{
  type=simbolos,
  name={$M_{res}$},
  description={Momento flector resultante en la sección crítica},
  unit={N$\cdot$m},
  sort={M-res}
}

\newglossaryentry{Dano}{
  type=simbolos,
  name={$D$},
  description={Daño acumulado por fatiga (regla de Palmgren-Miner); la verificación exige $D \leq 1$},
  unit={---},
  sort={D-1}
}

\newglossaryentry{D20}{
  type=simbolos,
  name={$D_{20}$},
  description={Daño acumulado al cabo de los $20$ años de vida de diseño},
  unit={---},
  sort={D-20}
}

\newglossaryentry{ni}{
  type=simbolos,
  name={$n_i$},
  description={Número de ciclos aplicados en el rango de amplitud $i$},
  unit={---},
  sort={n-i}
}

\newglossaryentry{Ni}{
  type=simbolos,
  name={$N_i$},
  description={Número de ciclos admisibles en el rango de amplitud $i$},
  unit={---},
  sort={N-i}
}

\newglossaryentry{Nf}{
  type=simbolos,
  name={$N_f$},
  description={Número de ciclos hasta la falla bajo amplitud constante},
  unit={---},
  sort={N-f}
}

\newglossaryentry{lambdaopt}{
  type=simbolos,
  name={$\lambda_{opt}$},
  description={Relación de velocidad de punta de máxima eficiencia},
  unit={---},
  sort={lambda-opt}
}

\newglossaryentry{Taire}{
  type=simbolos,
  name={$T$},
  description={Temperatura del aire en el emplazamiento},
  unit={K},
  sort={T-aire}
}

\newglossaryentry{mu0}{
  type=simbolos,
  name={$\mu_0$},
  description={Viscosidad de referencia de la ley de Sutherland},
  unit={Pa$\cdot$s},
  sort={mu-0}
}

\newglossaryentry{T0}{
  type=simbolos,
  name={$T_0$},
  description={Temperatura de referencia de la ley de Sutherland},
  unit={K},
  sort={T-0}
}

\newglossaryentry{Csuth}{
  type=simbolos,
  name={$C$},
  description={Constante de Sutherland para el aire},
  unit={K},
  sort={C-suth}
}

\newglossaryentry{Ra}{
  type=simbolos,
  name={$R_a$},
  description={Constante del gas para el aire seco},
  unit={J/(kg$\cdot$K)},
  sort={R-a}
}

\newglossaryentry{Patm}{
  type=simbolos,
  name={$P_{atm}$},
  description={Presión atmosférica},
  unit={Pa},
  sort={P-atm}
}
